\chapter{Návrh}
\label{chap:navrh}

Tato kapitola popisuje návrh agenta: jeho architekturu, numerickou doménu \texttt{sc2-terran-macro}, způsob, jakým doména modeluje čas, etapové cíle, vykonávací a~reaktivní vrstvu a~vývoj domény přes jednotlivé verze. Na konci je popsána klasická doména, se kterou se numerický agent srovnává. Popis odpovídá verzi kódu \texttt{a3cb452}, na které běží hlavní série experimentů (kapitola~\ref{chap:experimenty}), a~doméně \texttt{domain\_numeric\_v11.pddl}, jejíž úplné znění je v~příloze~\ref{app:domena}.

\section{Architektura agenta}
\label{sec:architektura}

Agent je rozdělen na plánovací vrstvu, která rozhoduje, vykonávací vrstvu, která rozhodnutí převádí na příkazy hry a~ověřuje jejich účinek, a~úzkou reaktivní vrstvu pro situace, které nemohou čekat na plánovač (obrázek~\ref{fig:architektura}). Řízení probíhá v~cyklu \emph{plán $\to$ zadání $\to$ ověření $\to$ přeplánování}:

\begin{enumerate}
	\item Třída \texttt{PlanningBot} určí cíl aktuální etapy (sekce~\ref{sec:etapy}).
	\item Modul \texttt{dynamic\_problem} přečte z~pozorování hry hodnoty všech fluentů a~zapíše je spolu s~cílem jako problém PDDL. Problém se generuje znovu před každým voláním plánovače, takže plánovač vždy vychází ze skutečného stavu.
	\item Modul \texttt{planner} spustí nad doménou a~problémem plánovač ENHSP s~časovým limitem 10~s a~souběžně problém načte knihovnou unified-planning, aby šlo plán dále ověřovat v~modelu. Výsledkem je posloupnost akcí bez parametrů, například \texttt{wait\_for\_minerals}, \texttt{build\_barracks}, \texttt{train\_marine}.
	\item Akce se vykonávají po jedné. Modul \texttt{executor} akci převede na příkaz rozhraní hry a~ohlásí, zda jej mohl zadat. Zadaná akce se považuje za hotovou teprve tehdy, když její \emph{podmínka dokončení} potvrdí účinek ve hře (sekce~\ref{sec:vykonavani}). Když hra akci zadržuje, smí se před ni posunout pozdější akce téhož plánu, pokud model potvrdí, že přeuspořádaný plán stále platí.
	\item Když se plán vyčerpá, agent postoupí do další etapy a~naplánuje znovu; když akci nelze zadat, přeplánuje aktuální etapu.
\end{enumerate}

\begin{figure}[ht]
	\centering
	% Agent architecture (chapter Návrh). Drawn in TikZ; module names match the repository.
% Needs \usetikzlibrary{babel,positioning,arrows.meta,calc} in the preamble.
\begin{tikzpicture}[
		node distance=11mm and 12mm,
		box/.style={draw, rounded corners=2pt, align=center, font=\scriptsize,
			minimum width=30mm, minimum height=9mm, inner sep=2pt},
		ext/.style={box, fill=black!8},
		flow/.style={-{Stealth[length=1.8mm]}, semithick},
		obs/.style={flow, dashed},
		lbl/.style={font=\scriptsize, align=center, fill=white, inner sep=1pt},
	]

	% columns: planning | problem and verification | planner and reactive layer
	\node[box] (bot) {\texttt{PlanningBot}\\ etapy, \texttt{on\_step}};
	\node[box, right=of bot] (problem) {\texttt{dynamic\_problem}\\ stav hry $\to$ PDDL};
	\node[box, right=of problem] (planner) {\texttt{planner}\\ unified-planning};

	\node[box, below=of bot] (exec) {\texttt{executor}\\ akce $\to$ příkaz hry};
	\node[box, below=of problem] (gate) {podmínky dokončení\\ vůči baseline};
	\node[ext, below=of planner] (enhsp) {ENHSP (JVM)\\ doména \texttt{v9}};

	\node[box, below=of exec] (combat) {\texttt{combat}, \texttt{placement}\\ útok, obrana, umístění};
	\node[box, below=of gate] (reactive) {reaktivní vrstva\\ supply, obrana, dělníci};
	\node[box, below=of enhsp] (metrics) {\texttt{metrics}\\ události, snímky};

	\node[ext, below=9mm of reactive, minimum width=104mm] (sc2) {StarCraft~II (burnysc2)};

	% planning
	\draw[flow] (bot) -- node[lbl, above=1pt] {cíl} (problem);
	\draw[flow] (problem) -- node[lbl, above=1pt] {problém} (planner);
	\draw[flow] ([xshift=-3mm]planner.south) -- ([xshift=-3mm]enhsp.north);
	\draw[flow] ([xshift=3mm]enhsp.north) -- ([xshift=3mm]planner.south);
	\draw[flow] (planner.north) to[bend right=14] node[lbl, above=1pt] {plán} (bot.north);

	% execution
	\draw[flow] (bot) -- node[lbl, left=1pt] {akce} (exec);
	\draw[flow] (exec) -- node[lbl, below=1pt] {zadáno} (gate);
	\draw[flow] (gate.north west) -- node[lbl, pos=0.5] {další akce /\\ přeplánovat} (bot.south east);
	\draw[flow] (exec) -- (combat);
	\draw[flow] (combat) -- node[lbl, left=1pt] {příkazy} (combat |- sc2.north);
	\draw[flow] (reactive) -- (sc2.north -| reactive);

	% observation
	\draw[obs] (sc2.north -| metrics) -- node[lbl, right=1pt] {stav} (metrics);
	\draw[obs] (sc2.west) -- ++(-5mm,0) |- node[lbl, pos=0.25, left=1pt] {stav hry} (bot.west);
\end{tikzpicture}

	\caption[Architektura agenta.]{Architektura agenta. Plné šipky jsou tok řízení, čárkované tok pozorování. Šedé bloky jsou externí systémy.}
	\label{fig:architektura}
\end{figure}

Oproti~\cite{Orava2025} se architektura liší ve dvou bodech. Plánovač dostává úplný numerický stav (suroviny, toky, kapacity), ne jen fakta o~existenci budov, a~akce se neodškrtává zadáním příkazu, ale až potvrzeným účinkem. Druhý bod je nutný právě kvůli prvnímu: numerický plán obsahuje předpoklady o~tom, kolik surovin v~kterém okamžiku bude, a~ty musí vykonávání dodržet, jinak se plán a~hra rozejdou (sekce~\ref{sec:limity}).

\section{Numerická doména}
\label{sec:domena}

Doména \texttt{sc2-terran-macro} je zapsána v~PDDL~2.1 úrovně~2 a~deklaruje jen požadavky \texttt{:typing} a~\texttt{:numeric-fluents}. Nemá objekty ani predikáty: veškerý stav je ve 24~numerických fluentech (tabulka~\ref{tab:fluenty}) a~akce jsou bez parametrů. Volba počtů místo objektů je záměrná. Agent nepotřebuje rozlišovat jednotlivé budovy téhož typu, protože výběr konkrétní budovy, na kterou se příkaz zadá, je úlohou vykonávací vrstvy, a~akce bez parametrů nemusí plánovač instancovat přes všechny objekty, takže groundovaný problém zůstává malý bez ohledu na velikost armády.

\begin{table}[ht]
	\centering
	\caption{Fluenty domény \texttt{domain\_numeric\_v11.pddl}.}
	\label{tab:fluenty}
	\begin{tabularx}{\textwidth}{@{} l >{\raggedright\arraybackslash}X @{}}
		\toprule
		\textbf{Skupina} & \textbf{Fluenty} \\
		\midrule
		Suroviny a~toky & \texttt{minerals}, \texttt{vespene}, \texttt{mining-rate}, \texttt{gas-rate} \\
		Supply & \texttt{supply-used}, \texttt{supply-cap} \\
		Budovy & \texttt{depot-count}, \texttt{barracks-count}, \texttt{refinery-count}, \texttt{factory-count}, \texttt{engineeringbay-count}, \texttt{cc-count}, \texttt{techlab-barracks-count}, \texttt{techlab-factory-count} \\
		Armáda & \texttt{marine-count}, \texttt{marauder-count}, \texttt{tank-count} \\
		Technologie & \texttt{weapon-upgrades}, \texttt{armor-upgrades}, \texttt{stimpack} \\
		Výrobní kapacita & \texttt{rax-time}, \texttt{factory-time} \\
		Útok & \texttt{attacks-launched}, \texttt{attack-threshold} \\
		\bottomrule
	\end{tabularx}
\end{table}

Doména má 16~akcí (tabulka~\ref{tab:akce}). Ceny všech akcí odpovídají cenám, které hra skutečně účtuje; kontroluje to automatický test, který z~domény vyčte efekty \texttt{decrease} nad fluenty \texttt{minerals} a~\texttt{vespene} a~porovná je s~tabulkou cen hry (\texttt{tests/test\_domain\_costs.py}). Test vznikl poté, co se ukázalo, že všechny verze až po čtvrtou účtovaly stavbu Factory bez jejích 100~jednotek plynu (sekce~\ref{sec:vyvoj}).

\begin{table}[ht]
	\centering
	\caption[Akce domény \texttt{domain\_numeric\_v11.pddl}.]{Akce domény \texttt{domain\_numeric\_v11.pddl}. Cena je v~minerálech/plynu; všechny akce kromě akcí čekání navíc vyžadují \texttt{(< (attacks-launched) 1)}.}
	\label{tab:akce}
	\footnotesize
	\begin{tabularx}{\textwidth}{@{} l r >{\raggedright\arraybackslash}X >{\raggedright\arraybackslash}X @{}}
		\toprule
		\textbf{Akce} & \textbf{Cena} & \textbf{Další předpoklady} & \textbf{Efekty mimo cenu} \\
		\midrule
		\texttt{wait\_for\_minerals} & -- & -- & $+$\texttt{mining-rate} minerálů, $+$\texttt{gas-rate} plynu, $+5$\,s výroby na budovu \\
		\texttt{wait\_for\_vespene} & -- & Refinery & totéž \\
		\texttt{build\_supply\_depot} & 100/0 & cap $\le 192$, volné supply $< 16$ & cap $+8$ \\
		\texttt{build\_barracks} & 150/0 & depot, Barracks $< 4\cdot$CC & Barracks $+1$ \\
		\texttt{build\_refinery} & 75/0 & Refinery $< 2\cdot$CC & \texttt{gas-rate} $+25$ \\
		\texttt{expand} & 400/0 & CC $< 3$ & CC $+1$, \texttt{mining-rate} $+40$ \\
		\texttt{build\_factory} & 150/100 & Barracks, Factory $<$ CC & Factory $+1$ \\
		\texttt{build\_engineeringbay} & 125/0 & depot, žádná EBay & EBay $+1$ \\
		\texttt{build\_techlab\_barracks} & 50/25 & Barracks bez add-onu & Tech Lab $+1$ \\
		\texttt{build\_techlab\_factory} & 50/25 & Factory bez add-onu & Tech Lab $+1$ \\
		\texttt{train\_marine} & 50/0 & Barracks, \texttt{rax-time} $\ge 18$, 1~supply & Marine $+1$ \\
		\texttt{train\_marauder} & 100/25 & Tech Lab, \texttt{rax-time} $\ge 21$, 2~supply & Marauder $+1$ \\
		\texttt{train\_siegetank} & 150/125 & Tech Lab u~Factory, \texttt{factory-time} $\ge 32$, 3~supply & Tank $+1$ \\
		\texttt{research\_weapons} & 100/100 & EBay, úroveň $< 1$ & úroveň $+1$ \\
		\texttt{research\_stimpack} & 100/100 & Tech Lab u~Barracks & Stimpack \\
		\texttt{attack} & -- & hodnota armády $\ge$ \texttt{attack-threshold} & armáda $:= 0$, \texttt{attacks-launched} $+1$ \\
		\bottomrule
	\end{tabularx}
\end{table}

Model zjednodušuje hru v~několika ohledech, které je třeba znát při čtení výsledků. Doména neobsahuje dělníky jako akci; rychlost těžby se před každým plánováním odvozuje z~jejich skutečného počtu a~nová základna nebo Refinery ji zvýší pevným odhadem (\texttt{mining-rate} $+40$, \texttt{gas-rate} $+25$). Expanze je omezena na tři základny. Engineering Bay vyžaduje v~modelu Supply Depot, ač ve hře stačí Command Center. Fluent \texttt{armor-upgrades} je deklarován, ale žádná akce ho nemění, a~výzkum zbraní je jen jedné úrovně, protože Armory potřebná pro další úrovně v~modelu není. Limity výroby (nejvýše čtyři Barracks a~jedna Factory na základnu, jedna Engineering Bay, depoty jen do kapacity 192 a~jen při méně než 16~volných supply) nejsou pravidla hry, ale ořez stavového prostoru: bez nich plánovač stavěl budovy, které pak neměl čím zaplnit, a~s~příliš těsnými limity zase neměl za co utratit nastřádanou banku (sekce~\ref{sec:vyvoj}).

Akce \texttt{attack} (výpis~\ref{lst:attack}) vyžaduje, aby vážená hodnota armády, Marine~1, Marauder~2 a~Siege Tank~3, dosáhla prahu \texttt{attack-threshold}, a~armádu \enquote{spotřebuje}: vynuluje počty jednotek a~zvýší počet vyslaných útoků. Protože všechny ostatní akce kromě čekání vyžadují \texttt{(< (attacks-launched) 1)}, je útok v~každém plánu nejvýše jeden a~je jeho poslední akcí; vše, co plán postaví a~vytrénuje, se odehraje před ním.

\begin{lstlisting}[caption={Akce útoku v~doméně v11.},label=lst:attack,float,basicstyle=\ttfamily\small]
(:action attack
  :parameters ()
  :precondition (and (< (attacks-launched) 1)
    (>= (+ (marine-count) (+ (* 2 (marauder-count))
                             (* 3 (tank-count))))
        (attack-threshold)))
  :effect (and (assign (marine-count) 0)
               (assign (marauder-count) 0)
               (assign (tank-count) 0)
               (increase (attacks-launched) 1)))
\end{lstlisting}

\section{Modelace času}
\label{sec:cas}

Doména neobsahuje durativní akce; čas do ní vstupuje dvěma akcemi čekání (výpis~\ref{lst:wait}). Jedna akce čekání představuje pevný úsek $\Delta = 5$~herních sekund (\texttt{WAIT\_SECONDS}) a~přičte k~fluentům to, co se za tento úsek nastřádá: minerály podle \texttt{mining-rate}, plyn podle \texttt{gas-rate} a~výrobní sekundy podle počtu výrobních budov. Plán je díky tomu nejen pořadím akcí, ale i~hrubým rozvrhem: počet akcí čekání mezi dvěma akcemi říká, kolik herního času mezi nimi uplyne, než si je agent může dovolit.

\begin{lstlisting}[caption={Akce čekání (\enquote{hodiny}) v~doméně v11.},label=lst:wait,float,basicstyle=\ttfamily\footnotesize]
(:action wait_for_minerals
  :parameters () :precondition ()
  :effect (and (increase (minerals) (mining-rate))
               (increase (vespene) (gas-rate))
               (increase (rax-time) (* 5 (barracks-count)))
               (increase (factory-time) (* 5 (factory-count)))))
\end{lstlisting}

\paragraph{Rychlost těžby.} Hodnoty \texttt{mining-rate} a~\texttt{gas-rate} jsou množství suroviny, které hra skutečně vytěží za $\Delta$, a~počítají se z~dělníků právě přiřazených k~těžbě:
\[
	\texttt{mining-rate} = \left\lfloor \Delta \cdot r_m \cdot \sum_{b} \min(h_b, 16) \right\rfloor,
	\qquad
	\texttt{gas-rate} = \left\lfloor \Delta \cdot r_g \cdot \sum_{r} h_r \right\rfloor,
\]
kde $h_b$ je počet dělníků na minerálech u~základny~$b$, $h_r$ počet dělníků u~Refinery~$r$ a~$r_m = 0{,}93$ a~$r_g = 0{,}84$ jsou rychlosti těžby jednoho dělníka za sekundu. Hodnota~$r_g$ byla kalibrována podle her sérií T8.1 a~T8.2 a~platí od verze~7 domény (\serie{T8.2_v7_medium}). Při dvanácti dělnících na začátku hry vychází 55~minerálů na jedno čekání. Obě hodnoty jsou nejméně~1, protože akce s~nulovým přírůstkem nic nemění a~plánovač by problém, ve kterém je potřeba čekat, prohlásil za neřešitelný. Stejnou funkci používá i~vykonávací vrstva, když rozhoduje, zda čekání skončilo, takže to, co předpokládá doména, a~to, co se ve hře testuje, se nemůže rozejít.

\paragraph{Výrobní kapacita.} Od verze~6 nese doména dva fluenty výrobního času, \texttt{rax-time} a~\texttt{factory-time}: kolik sekund výroby mají Barracks, resp. Factory k~dispozici. Každé čekání přidá 5~s za každou výrobní budovu a~trénink jednotky spotřebuje její dobu výroby (Marine~18~s, Marauder~21~s, Siege Tank~32~s). Na začátku plánu se fluenty nastaví podle budov, které jsou právě volné. Plán, který chce utratit velkou banku rychle, tak musí postavit víc výrobních budov; jinak by mohl naplánovat dvacet tanků z~jedné Factory, které by se ve hře vyráběly déle než deset minut. Model je přitom hrubý: výrobní čas \texttt{rax-time} roste o~5~s za \emph{každé} Barracks, i~když Marauder potřebuje Barracks s~Tech Labem. Plánovač proto může počítat s~Maraudery ze všech Barracks, zatímco hra je vyrábí jen u~jednoho; v~hlavní sérii proti VeryEasy to ve dvou hrách vedlo k~plánům s~25 a~28~Maraudery, které držely armádu doma přes deset minut (sekce~\ref{sec:vs-ai}).
% TODO (T4.9): doménu v12 s výrobním časem Tech Labů doplnit, až bude změřena.

\paragraph{Jedny hodiny.} Ve verzi~6 přičítala akce \texttt{wait\_for\_vespene} jen plyn a~výrobní čas, nikoli minerály vytěžené za stejnou dobu, a~plánovač ji používal jako levný zdroj výrobního času. Od verze~7 přičítají obě akce čekání všechny tři veličiny současně; liší se jen tím, že \texttt{wait\_for\_vespene} vyžaduje Refinery. Model má tedy jedny hodiny a~každé čekání je pět sekund pro celou ekonomiku.

\paragraph{Proč ne durativní akce.} Spojitý model těžby, buď durativní akcí se spojitým efektem \texttt{(increase (minerals) (* \#t (mining-rate)))}, nebo procesem PDDL+~\cite{FoxLong2006}, by byl věrnější: plán by byl rozvrhem s~makespanem a~čekání by nebylo třeba plánovat. Diskrétní hodiny byly zvoleny ze tří důvodů. Plán zůstává posloupností okamžitých akcí, takže vykonávání i~podmínky dokončení jsou jednoduché a~každou akci čekání lze ve hře přímo ověřit. Na doménu stačí plánovač úrovně~2, kterých je víc a~jsou rychlejší než plánovače se spojitým časem. A~rozlišení pěti sekund je vůči době výroby jednotek (18 až 32~s) a~staveb (desítky sekund) dostatečně jemné. Cenou je délka plánů: akce čekání tvoří velkou část každého plánu (na pevných problémech etap 0 až 2 jde o~18 z~28, 14 z~26 a~18 z~36 akcí, \serie{T5.4_planners_offline}) a~vyhovující plánovač bez metriky někdy čeká víc, než je nutné. Akce čekání přitom nelze z~domény odstranit: bez nich by čísla v~doméně jen omezovala pořadí akcí a~numerický model by se zredukoval na klasický.

Spojitou variantu práce ověřila mimo hru. Doména \texttt{domain\_temporal.pddl} je doména v11 zapsaná v~PDDL+: akce čekání nahrazuje proces, který po celou dobu spojitě zvyšuje suroviny a~výrobní čas, a~ENHSP ji plánuje s~rozhodovacím krokem 1~s. Na pevných problémech etap 0 až 2 našel plány za přibližně 0,6~s s~délkou rozvrhu (\emph{makespan}) 76, 73 a~122~herních sekund, zatímco diskrétní plány téhož plánovače odpovídají 90, 70 a~90~sekundám (\serie{T5.4_planners_offline}). Spojitý model je tedy řešitelný stejně rychle; v~etapě~0 dává kratší plán, v~etapě~2 delší, protože obě varianty hledají vyhovující, nikoli nejkratší plán. Ve hře se spojitá varianta nepoužívá: plán s~časovými značkami by vyžadoval jiný druh vykonávání, který akce nespouští po jedné, ale podle času, a~ten práce neimplementuje. Ani jedna z~variant zatím nemodeluje dobu staveb a~výzkumů; ve hře se to projevuje například tím, že depot přidá supply v~modelu okamžitě, ve hře až po 21~s stavby (\serie{T4.8_supply_wait_veryeasy}). Spojitý čas by takové trvání vyjádřit uměl, což je přirozené pokračování práce.

\section{Etapové cíle a~generování problému}
\label{sec:etapy}

Hra je rozdělena do čtyř etap s~cíli v~tabulce~\ref{tab:etapy}. Etapy 0 až 2 budují technologii a~armádu po krocích, které odpovídají technologickému stromu; etapa~3 je poslední a~opakuje se. Její cíl \texttt{(>= (attacks-launched) 1)} je splněn akcí \texttt{attack} a~generátor problému nastavuje \texttt{attacks-launched} vždy na~0, takže cíl znamená \enquote{vyšli ještě jednu vlnu}. Aby plánovač mohl zaútočit, musí nejprve postavit armádu s~hodnotou alespoň \texttt{attack-threshold}; po vyčerpání plánu se etapa~3 naplánuje znovu. Pozdní hra je tak cyklem shromáždění, útoku a~nového shromáždění, o~jehož načasování i~složení rozhoduje plánovač.

\begin{table}[ht]
	\centering
	\caption{Etapové cíle (\texttt{STAGE\_GOALS}).}
	\label{tab:etapy}
	\begin{tabularx}{\textwidth}{@{} c >{\raggedright\arraybackslash}X @{}}
		\toprule
		\textbf{Etapa} & \textbf{Cíl} \\
		\midrule
		0 & $\ge 1$ Barracks, $\ge 6$ Marine \\
		1 & $\ge 1$ Refinery, $\ge 1$ Tech Lab u~Barracks, $\ge 4$ Marauder, $\ge 8$ Marine \\
		2 & $\ge 1$ Factory, $\ge 1$ Tech Lab u~Factory, $\ge 2$ Siege Tank, zbraně úrovně~1, Stimpack, $\ge 14$ Marine \\
		3 (opakuje se) & $\ge 1$ vyslaný útok \\
		\bottomrule
	\end{tabularx}
\end{table}

\paragraph{Strop banky.} Samotné cíle etap neříkají nic o~surovinách, které agent nastřádal navíc. Kdyby plánovač jen splnil cíl, banka by rostla, což je stagnace popsaná v~\cite{Orava2025}. Každý cíl je proto rozšířen o~podmínku na zůstatek minerálů na konci plánu:
\[
	\textit{cíl etapy} \wedge \big(\texttt{supply-used} \ge 190 \vee \texttt{minerals} \le c\big).
\]
Agent nejprve plánuje s~přísným stropem $c = 400$, takže plán musí utratit celou banku nad tento zůstatek a~rozhodnout, za co. Když pro přísný cíl plán nenajde, zkusí slabší strop $c = \max(400,\; m_0 - 1\,500)$, kde $m_0$ je stav banky při plánování, tedy utratit nejvýše 1\,500~minerálů. Do verze kódu, na které běžely série T8, platil jen slabší strop; část příjmu jednoho plánovacího cyklu pak v~bance zůstávala (\serie{T4.8_spend_bank_veryeasy}). Podmínka odpadá, když už se do supply nevejdou další jednotky. Na plyn strop není: s~několika Refinery přinese každé čekání zhruba tolik plynu, kolik jednotka stojí, takže plynový strop byl nedosažitelný a~prodlužoval plány na stovky akcí (\serie{T8.3_v8_medium}).

\paragraph{Práh útoku.} Útok vyžaduje armádu s~hodnotou 30. Jednotky vyslané do útoku ale dál drží supply, a~když armáda v~poli zabírá většinu kapacity, doma se armáda s~hodnotou 30 postavit nedá a~plánovač by marně hledal plán, dokud nevyprší limit. Generátor problému proto nastavuje
\[
	\texttt{attack-threshold} = \max\big(0, \min(30,\; v_{\mathrm{doma}} + c_{\mathrm{dosažitelná}} - \texttt{supply-used})\big),
\]
kde $v_{\mathrm{doma}}$ je hodnota armády u~základen a~$c_{\mathrm{dosažitelná}}$ kapacita supply, kterou lze ještě postavit depoty (193 až 200 podle aktuální kapacity). Když je práh nulový, nevejde se doma nic a~cílem etapy~3 je jen jedno čekání.

Problém PDDL vzniká přímo z~pozorování: každý fluent je funkcí stavu hry (počty budov a~jednotek, banka, rychlosti těžby, volné výrobní budovy). Armáda se počítá jen z~jednotek, které jsou u~základen a~nebyly vyslány do probíhajícího útoku; ty už plánovač nevidí (sekce~\ref{sec:vykonavani}). Generátor zapisuje jen fluenty, které daná doména deklaruje, takže agent může hrát s~libovolnou starší verzí domény, což je podmínkou reprodukovatelnosti historických sérií.

\section{Vykonávací vrstva}
\label{sec:vykonavani}

Vykonávací vrstva odděluje dvě otázky: zda lze akci \emph{zadat} a~zda se akce \emph{dokončila}. Na první odpovídá \texttt{executor}: vrátí, zda příkaz zadal, a~pokud ne, uvede důvod (například \enquote{cannot afford}, \enquote{no free geyser}). Důvod začínající slovem \texttt{waiting} znamená, že akce jen čeká na dokončení rozestavěné budovy nebo na uvolnění fronty, což je postup, nikoli selhání.

Na druhou otázku odpovídá \emph{podmínka dokončení} (tabulka~\ref{tab:podminky}). Podmínky se vyhodnocují proti \emph{baseline}, stavu zachycenému těsně před zadáním akce. Díky tomu plán se čtyřmi akcemi \texttt{build\_supply\_depot} postaví čtyři depoty, a~ne jeden, který by splnil podmínku \enquote{nějaký depot existuje}, a~čekání skončí, až banka vzroste o~\texttt{mining-rate} proti stavu při zadání. Vykonávací vrstva tedy nikdy nepostoupí na základě předpokladu, že akce vyšla; vždy čeká, až to potvrdí hra.

\begin{table}[ht]
	\centering
	\caption{Podmínky dokončení akcí vůči baseline zachycené při zadání.}
	\label{tab:podminky}
	\begin{tabularx}{\textwidth}{@{} l >{\raggedright\arraybackslash}X @{}}
		\toprule
		\textbf{Akce} & \textbf{Akce je hotová, když} \\
		\midrule
		stavba, expanze, add-on & počet budov daného typu (hotových i~rozestavěných) přesáhl baseline \\
		trénink & počet jednotek daného typu (živých i~ve frontě) přesáhl baseline \\
		výzkum & vylepšení je dokončené nebo prokazatelně probíhá \\
		čekání & banka dosáhla stavu při zadání $+$ \texttt{mining-rate} (resp.\ \texttt{gas-rate}) \\
		útok & žádná z~vyslaných jednotek nestojí nečinně u~základen \\
		\bottomrule
	\end{tabularx}
\end{table}

\paragraph{Časové limity a~eskalace.} Zadaná akce, která se do 30~s nedokončí (například dělník zahynul cestou na staveniště), se uvolní a~zaznamená jako \texttt{action\_timeout}. Akci, kterou nelze zadat 30~s, řeší agent eskalací ve třech krocích: nejprve etapu přeplánuje z~aktuálního stavu, takže plánovač sám uvidí, co chybí, a~buď to postaví, nebo cíl obejde. Po třech takových přeplánováních bez mezilehlého postupu se další nevykonatelné akce vynechávají a~po dalších třech agent opustí celou etapu. Akce, která čeká jen na dokončení stavby, smí čekat až 120~s. Počítadlo pokusů se nuluje jen skutečným postupem, tedy zadáním nebo dokončením akce; nulování při vynechání dovolilo agentovi v~jedné z~dřívějších verzí donekonečna cyklit mezi dvěma nevykonatelnými akcemi.

\paragraph{Přeuspořádání plánu.} Plán se vykonává po jedné akci, ale hra některou akci často zadrží: Barracks pro Tech Lab se ještě staví, všechny výrobní fronty jsou plné, jednotka čeká na rozestavěný depot. Dokud se taková akce nedala zadat, stál celý plán; v~sérii proti VeryEasy na doméně v9 to bylo přibližně 290~s z~přibližně 810~s hry a~příjem těchto sekund model neviděl (\serie{T5.2_VeryEasy}). Když vykonávací vrstva hlásí, že akce jen čeká, posune agent před ni první pozdější akci téhož plánu, ale jen tehdy, když přeuspořádaný plán z~počátečního stavu problému podle simulátoru knihovny unified-planning stále projde všemi předpoklady a~skončí v~cíli (modul \texttt{plan\_order}). Akce čekání se nepřesouvají. Agent tak nerozhoduje o~ničem novém: mění se jen pořadí akcí, které vybral plánovač, a~jen tam, kde to dovoluje jeho vlastní model. Každý posun se zaznamená jako \texttt{plan\_reordered}.

\paragraph{Selhání plánovače.} Agent zkouší cíle od nejsilnějšího: nejprve cíl s~přísným stropem banky, potom se slabším (sekce~\ref{sec:etapy}). Pro každý cíl spustí ENHSP ve výchozí konfiguraci ($h_{\mathrm{add}}$, GBFS) a~při vypršení časového limitu ještě jednou s~konfigurací $h_{\mathrm{mrp}}$ a~GBFS; během vývoje se ukázalo, že obě konfigurace selhávají na různých stavech. Teprve když plán nenajde žádná kombinace, přejde agent do nouzového režimu: trénuje Marine a~při šesti Marinech zaútočí. Nouzový režim je jediným místem, kde o~útoku nerozhoduje plánovač, a~každý neúspěšný pokus se zaznamená jako \texttt{planner\_error}.

\paragraph{Útok mimo plán.} Akce \texttt{attack} v~doméně armádu spotřebuje a~vykonávání to dodržuje doslova. Vykonávací vrstva útok zadá, až když ve hře platí jeho předpoklad z~modelu: hodnota armády u~základen dosáhla prahu, se kterým plán počítal. Trénink se totiž považuje za hotový už zařazením do fronty, takže část jednotek, se kterými plán počítal, ještě nemusí existovat; jednotky, které jsou při útoku stále ve výrobě, zůstanou doma pro další vlnu. Útok je hotový, jakmile všechny vyslané jednotky opustily základny; jejich identifikátory se uloží jako \emph{nasazené} a~plánovač je v~dalších problémech nevidí. Agent tak hned začne shromažďovat další vlnu, zatímco nasazené jednotky bojují mimo plán: postupují k~nejbližší známé budově soupeře, a~když žádnou neznají, prohledávají místa možných expanzí. Jakmile nasazené jednotky ztratí polovinu vyslané hodnoty, přeživší se vrátí domů a~připojí se k~další vlně. Rozhodnutí \emph{kdy} a~\emph{s~čím} zaútočit tedy patří plánovači; průběh boje je provedením tohoto rozhodnutí.

\section{Reaktivní vrstva}
\label{sec:reaktivni}

V~každém kroku hry před plánovací vrstvou běží čtyři reaktivní rutiny:
\begin{itemize}
	\item \textbf{rozdělování dělníků} na minerály a~plyn vestavěnou funkcí knihovny,
	\item \textbf{trénink dělníků} do saturace: 16 na Command Center a~3 na Refinery, nejvýše 66,
	\item \textbf{nouzový Supply Depot}, když zbývá méně než 2~volné supply a~žádný depot se nestaví,
	\item \textbf{obrana základny}: nepřátelské jednotky v~okruhu 35 od kterékoli vlastní základny napadne celá armáda, která je v~tom okruhu (jednotky útočící jinde se nevolají zpět); Siege Tanky se v~dosahu nepřítele rozloží.
\end{itemize}

Každá z~rutin řeší situaci, na kterou nelze čekat do dalšího plánu (sekce~\ref{sec:pozadavky}). Nouzový depot a~obrana se zaznamenávají jako \texttt{reactive\_override}, obrana jednou za střet, ne za každý krok hry, takže jejich rozsah je měřitelný: v~hlavní sérii proti obtížnosti Medium připadalo na hru v~průměru 219,8~zadaných akcí plánu a~24,1 takových zásahů (\serie{E5_numeric_Medium}). Rozdělování a~trénink dělníků se nezaznamenávají a~v~tomto poměru nejsou. Trénink dělníků je ze čtyř rutin nejméně obhajitelný, protože jde o~ekonomické rozhodnutí; jeho přesun do domény je otevřeným bodem (kapitola~\ref{chap:experimenty}).

\section{Vývoj domény v1--v11}
\label{sec:vyvoj}

Doména vznikala iterativně. Každá verze je samostatný soubor v~adresáři \texttt{planning/} (od \texttt{domain\_numeric.pddl} po \texttt{domain\_numeric\_v11.pddl}), vznikla jako odpověď na konkrétní pozorování z~her a~byla ověřena další sérií her. Tabulka~\ref{tab:vyvoj} shrnuje změny domény; změny vykonávací vrstvy v~rámci téže verze a~podrobný rozbor sérií proti obtížnosti Medium jsou v~kapitole~\ref{chap:experimenty}. Všechna čísla pocházejí z~uvedených sérií a~z~deníku experimentů \texttt{experiments/LOG.md}.

\begin{table}[p]
	\centering
	\caption[Vývoj numerické domény.]{Vývoj numerické domény. V~=~výhra, P~=~prohra, R~=~remíza (limit 1\,800 herních sekund); CI je 95\% Wilsonův interval.}
	\label{tab:vyvoj}
	\footnotesize
	\begin{tabularx}{\textwidth}{@{} c >{\raggedright\arraybackslash}X >{\raggedright\arraybackslash}X >{\raggedright\arraybackslash}p{3.6cm} @{}}
		\toprule
		\textbf{Verze} & \textbf{Změna} & \textbf{Výsledek} & \textbf{Série} \\
		\midrule
		v1 & 21 fluentů, 16 akcí; suroviny, supply, rychlosti těžby, počty budov a~jednotek & po opravách vykonávání výhry proti VeryEasy, v~etapě~3 se ale hromadí banka (4\,625 minerálů, 1\,344 plynu na konci hry) & \serie{2026-10-25_8d754f4}, \serie{2026-10-25_d8a042d} \\
		v2 & limity výroby, strop banky & etapa~3 vyrábí (Siege Tanky 0 $\to$ 12), banka ale roste dál (3\,975 $\to$ 5\,005 minerálů): jedna Factory nestačí, model nemá výrobní kapacitu & \serie{2026-10-25_e002cbb} \\
		v3 & zbraně jen 1.~úroveň, depot jen při $< 16$ volných supply; útok uvolňuje supply armády & hypotéza vyvrácena: model uvolnil supply, které přeživší armáda dál drží, a~jednotky byly odmítnuty při 77/79 supply & \serie{2026-10-26_8f239e5} \\
		v4 & útok supply neuvolňuje & výhry proti VeryEasy; porážka, když \texttt{build\_factory} odmítnuto při 3\,000 minerálech: doména neúčtovala plyn Factory & \serie{2026-10-26_bbf1fa4}, \serie{2026-10-26_06116be} \\
		v5 & Factory stojí 100 plynu; test všech cen & VeryEasy a~Easy 6/6~V; Medium 3~V / 6~P / 1~R, 30~\% (CI 11--60~\%); banka nad 500 minerály 87--95~\% hry & \serie{2026-10-30_T1.6_85bd1dc}, \serie{T8.1_medium} \\
		v6 & výrobní kapacita (\texttt{rax-time}, \texttt{factory-time}) & VeryEasy 3/3~V; recenze odhalila plynové čekání jako volný výrobní čas & \serie{T8.2_veryeasy}, \serie{T8.2_medium} \\
		v7 & jedny hodiny pro obě čekání, $r_g = 0{,}84$ & Medium 8~V / 2~P, 80~\% (CI 49--94~\%); medián délky hry 690~s (v5: 1\,063~s) & \serie{T8.2_v7_veryeasy}, \serie{T8.2_v7_medium} \\
		v8 & práh útoku 30 (dříve 20), útok končí plán & Medium 8~V / 0~P / 2~R, 80~\% (CI 49--94~\%), bez porážek & \serie{T8.3_v8_medium} \\
		v9 & práh útoku jako fluent omezený dosažitelným supply & Medium 4~V / 1~P, 80~\% (CI 38--96~\%) & \serie{T8.3_v9_medium}, \serie{T8.3_reachable_cap_medium} \\
		v10 & v9 $+$ \texttt{total-cost} pro pokusy s~metrikou plánu & metrika nezlepšila žádný plán; jen pro experimenty mimo hru & \serie{T4.2_planner_configs} \\
		v11 & Barracks $\le 4\cdot$CC, až 3~CC; plány utrácejí banku do 400 & VeryEasy 20/20~V; Medium 20/20 dokončených her~V, 3~starty zastaveny (sekce~\ref{sec:vs-ai}) & \serie{T4.8_spend_bank_veryeasy}, \serie{E5_numeric_VeryEasy}, \serie{E5_numeric_Medium} \\
		\bottomrule
	\end{tabularx}
\end{table}

Z~vývoje plynou dvě obecnější zjištění. Za prvé, většina verzí neopravovala plánovač, ale \emph{model}: chybějící cenu (v4 $\to$ v5), chybějící výrobní kapacitu (v5 $\to$ v6), nekonzistentní hodiny (v6 $\to$ v7). Teprve s~realistickým modelem kapacity a~času začal agent proti Medium vyhrávat většinu her. Za druhé, chyby modelu se projevovaly jako ekonomické symptomy, nejčastěji jako rostoucí banka, a~ty šly odhalit jen proto, že se ekonomika měřila (podíl hry s~bankou nad 500~minerály, hodnota armády v~čase). Stejný symptom, který~\cite{Orava2025} pozoruje jako stagnaci, je v~numerickém modelu diagnostickým signálem, který ukazuje na konkrétní chybějící část domény.

\section{Klasická doména pro ablaci}
\label{sec:klasicka}

Klasická doména \texttt{domain\_classical.pddl} slouží k~odpovědi na otázku, co přinášejí samotná čísla. Obsahuje stejné makroakce jako numerická doména kromě akcí čekání a~je zapsána v~propozičním stylu Oravovy domény~\cite{Orava2025}: každá budova a~jednotka je objekt, technologické předpoklady jsou existenční kvantifikátory a~stav jsou predikáty (\texttt{built}, \texttt{has-techlab}, \texttt{trained}, \texttt{weapons-researched}, \texttt{stim-researched}, \texttt{attacked}). Supply je zásobník slotů po dvou supply: každý depot \enquote{odemkne} čtyři sloty a~jednotka obsadí jeden slot bez ohledu na velikost. Suroviny, rychlosti těžby a~výrobní kapacita v~doméně nejsou a~cenu akcí hlídá až vykonávací vrstva, která příkaz odmítne, dokud na něj banka nestačí.

Generátor problému \texttt{classical\_problem.py} převádí stav hry na objekty a~fakta a~numerické etapové cíle na objekty, které je třeba postavit nebo vytrénovat; limity výroby verze~2 jsou vyjádřeny počtem deklarovaných objektů budov. Útočná vlna etapy~3 má pevné složení s~hodnotou 30 (12~Marine, 6~Marauder, 2~Siege Tanky), zmenšené stejně jako práh numerické domény, když se doma nevejde; vykonávací vrstva pak útok zadá při hodnotě této vlny. Plánovačem je Fast Downward v~konfiguraci LAMA-first, stejně jako v~\cite{Orava2025}. Vykonávací vrstva, podmínky dokončení, reaktivní vrstva i~metriky jsou sdílené, takže srovnání numerického a~klasického agenta se liší jen modelem a~plánovačem. Klasický agent je funkční: proti úrovni VeryEasy vyhrál kontrolní hru v~čase 744~s s~mediánem 13~akcí na plán (\serie{T4.1_classical_check}); systematické srovnání je v~kapitole~\ref{chap:experimenty}.
